%======================================================================================
\section{Cohesion: what keeps two pilots together}\label{sec:cohesion}

Inside a contact, the delay $\ell^*(t)$ says how many seconds one pilot is behind the other on the same
path. Instants with either pilot in a search are left out, as in $\ell_{\rm med}$: the pilots turn without
climbing and $\ell^*$ is not determined. Groups: the \StCohReal{} pair intervals of
pilots who took off within 10~km. Chance: the \StCohChance{} contacts of the surrogate pairs of the null
model (the same pairs as in the note), two pilots who never saw each other; the \StCohChanceLong{} that
last at least $\tau^*$ are selected as the groups are, by staying in contact that long.
\begin{itemize}
  \item $D(s)=\langle(\ell(t+s)-\ell(t))^2\rangle$, the mean square change of the delay after $s$, fitted by an
        Ornstein--Uhlenbeck process (Appendix~\ref{app:models}), $D(s)=2\sigma^2(1-e^{-ks})+c$: $1/k$ is how long
        the delay remembers its value, $\sigma$ its spread, $2k\sigma^2$ how fast it wanders.
  \item The drift $\langle\ell(t+60~\mathrm{s})-\ell(t)\rangle/60~\mathrm{s}$ at each $\ell$, made symmetric since the order of a pair is
        arbitrary: the pull back to $\ell=0$, a spring $-k\ell$ near 0.
  \item The distance between the two pilots at the same instant after their first contact.
\end{itemize}
Two independent pilots separate as twice the displacement variance of one, not twice its MSD: they share
the drift along the route. The surrogate pairs give this baseline directly, with the same terrain.

\begin{figure}[H]
\centering
\includegraphics[width=\textwidth]{cohesion.pdf}
\caption{Left: $D(s)$ of groups (by size) and of chance contacts, with $1/k$ of the Ornstein--Uhlenbeck fit.
Middle: symmetric drift of the delay. Right: median distance at the same instant after the first contact,
for every pair and for those whose first contact lasts at least $\tau^*$. Bands: 5--95\% of 200 draws of
the dates.}
\label{fig:cohesion}
\end{figure}

\begin{reading}
\textbf{Reading.}
\begin{itemize}
  \item \textbf{Groups keep the same pace.} Their delay wanders at $2k\sigma^2=\StQReal$ \StQRealBand~s$^2$/s, against
        \StQChanceLong{} \StQChanceLongBand{} for chance contacts as long and \StQChance{} \StQChanceBand{} for all of them; its
        spread is $\sigma=\StSigmaReal$~s \StSigmaRealBand, against \StSigmaChance~s \StSigmaChanceBand.
  \item \textbf{They are not pulled back much harder.} The delay forgets itself in $1/k=\StOuReal$~min \StOuRealBand, against
        \StOuChance~min \StOuChanceBand{} by chance; the drift gives \StDriftReal~min \StDriftRealBand{} against \StDriftChance~min
        \StDriftChanceBand. Most of the pull is the selection: a contact that lasts keeps $|\ell|\le\delta$ by definition.
  \item \textbf{Size does not matter:} $2k\sigma^2=\StQTwo$, \StQThree{} and \StQFour~s$^2$/s at 2, 3 and 4+ pilots.
  \item \textbf{They stay together after the contact that selected them.} When the first contact lasts at least $\tau^*$,
        the two pilots are \StSepRealLongSixty~km apart 60~min later (median), against \StSepChanceLongSixty~km by chance;
        \StSepRealLongTwoHours{} against \StSepChanceLongTwoHours~km after 2~h. Over every first contact, \StSepRealThirty{} against
        \StSepChanceThirty~km after 30~min.
\end{itemize}
A group is two pilots who keep the same pace, not two pilots tied by a spring. Whether this is following
or the same choices in the same air, the data cannot tell.
\end{reading}
